\chapter{Background}
\label{ch:background}

% Polymarket: spostato dal vecchio capitolo e completato; i concetti di ML sono da scrivere
This chapter gives the background that the rest of the thesis assumes: what a prediction
market is and how Polymarket works (\cref{sec:pm:what,ch:pm}), and the notions of machine
learning that the methodology uses: time series and sequence models (\cref{sec:bg:ts}),
graphs and graph neural networks (\cref{sec:bg:graphs}), learners on tables of features
(\cref{sec:bg:tabular}) and evaluation under class imbalance (\cref{sec:bg:eval}).

\section{Prediction markets}
\label{sec:pm:what}

A prediction market trades contracts whose value is tied to the outcome of a future
event.  The simplest contract is the \emph{binary} one: it pays one unit if the event
happens and nothing otherwise.  If traders are risk-neutral and have no budget
constraints, the price of such a contract equals the market's aggregate probability of
the event; under weaker assumptions the price is still close to the mean belief of the
traders \cite{wolfers2004prediction,manski2006interpreting}.  Surveys of the empirical
literature find that market prices are well calibrated and competitive with polls and
expert forecasts, and that they react to news within minutes
\cite{wolfers2004prediction,snowberg2012prediction,arrow2008promise}.  Two
systematic deviations are documented: a favourite--longshot bias, where unlikely events
are over-priced and likely events under-priced, which grows with the time to expiration
\cite{page2013calibrated}; and, on Polymarket itself, the same bias
\cite{cardozo2026favorite} together with a concentration of trading profits in a very
small set of accounts \cite{akey2026whowins}.

\section{Polymarket}
\label{ch:pm}

\subsection{How Polymarket works}
\label{sec:pm:mechanics}

Polymarket is a decentralised prediction market built on the Polygon blockchain, with
the stablecoin USDC as collateral.  Every market is a question with two outcomes:
``Yes'' and ``No'' for a binary question, ``Up'' and ``Down'' for a price interval, the
two teams of a game, ``Over'' and ``Under'' for a total.  Each outcome is a conditional
token: one dollar of collateral can at any time be \emph{split} into one Yes token and
one No token, a pair can be \emph{merged} back into one dollar, and after resolution the
winning token is \emph{redeemed} for one dollar and the losing one for nothing.  The two
prices therefore sum to about one dollar, and a trader who wants to sell Yes at price $p$
is indistinguishable, for the matching engine, from one who wants to buy No at $1 - p$.

Trading takes place on a central limit order book that is \emph{hybrid}.  Orders are
signed messages that the trader submits to an off-chain operator, which keeps the book,
publishes every change to it and matches crossing orders; a match is then settled
on-chain by the exchange contract, which transfers tokens and collateral and emits a fill
event with the maker, the taker, the token, the side, the quantity, the price and the
fee.  The book is thus observable as a stream of price-level updates, and the trades as
on-chain records attributed to wallet addresses.  The two streams are not the same
thing: most book events are orders placed or cancelled away from the best quotes and never
trade; a fill is a book event that found a counterparty.

Multi-outcome events, the candidates of an election or the strikes of a ``Bitcoin above
$x$'' ladder, are sets of binary markets whose prices are linked by the payoff structure
(exactly one candidate wins), with a ``negative risk'' conversion that lets a trader hold
the No side of every candidate as one position.  The prices of such siblings must respect
bounds that the market does not always enforce \cite{gebele2026arbitrage}.

\subsection{Tokens, contracts and transactions: the blockchain layer}
\label{sec:pm:chain}

Polygon produces a block roughly every two seconds; every fill is a transaction in one
of its blocks and carries the block number, the transaction hash and the wallet
addresses of the maker and the taker.  Because settlement is on-chain, every trade is
attributed to a wallet that persists across markets and across time, which makes it
possible to build a bipartite graph between wallets and markets, and to ask whether the
history of a wallet (its activity, its past accuracy, the other markets it trades) is
informative about the price of the markets it trades now.  Works on Polymarket have used
wallet-level data to study who wins and who loses \cite{akey2026whowins}, insider trading
and informed flow \cite{bsic2026polymarket}, and arbitrage between related markets
\cite{saguillo2025unravelling}; none of them, to our knowledge, has used the wallet--market
graph as the input of a forecasting model, which is one of the questions of this thesis
(\cref{ch:graph}).  The block is also the natural time step of the on-chain stream, and it
is the time step of every stateful model in the thesis (\cref{sec:data:blocks}); the graph
itself is described in \cref{sec:expl:graph}.

\subsection{The limit order book}
\label{sec:pm:lob}

A limit order book \cite{gould2013limit} is, for one token, two sorted lists of resting
limit orders, bids and asks, each with the total size resting at that price.  The highest
bid and the lowest ask are the best quotes, their average is the \emph{mid price} and
their difference the \emph{spread}; the size at the first levels on each side is the
\emph{depth}.  Prices on Polymarket live in $[0, 1]$ with a tick of $0.001$ on most
markets, so a contract that is almost certain trades at $0.99$ or above and its price can
only move by a few ticks.  Because the Yes and No books mirror each other, the book of one
side describes the market, and this thesis reads the Yes side.  An \emph{order-book event}
is a change of the size resting at one price on one side; most events do not touch the
best quotes.  A \emph{snapshot} is the full state of the book at one moment.

Two consequences of this structure recur in every result of the thesis.  A mid price is a
bounded, discrete process observed through a thin book, so most of its changes are
one-tick oscillations produced by an order arriving at or leaving the best quote, which
are not information and which revert; and the dynamics depend on the price level, since a
market at $0.98$ has two ticks of room on one side and a hundred on the other.
\textcite{dubach2026anatomy} document the resulting microstructure on a 52-day sample:
depth depends on the market's duration, price level and volume and is spread over hundreds
of levels rather than concentrated at the top, and the trade direction inferred from the
public book matches the on-chain record only about 60\% of the time.  This is why the
thesis takes the direction of a trade from the on-chain fill and not from the book.

\subsection{The life of a market}
\label{sec:pm:life}

A market is created with its question, its resolution rules, its two outcomes and an end
date; its contract is deployed on-chain and it starts accepting orders at known times, both
recorded in the market's metadata and both known in advance.  It then trades until the
event is decided.  Some markets are born and die within minutes: the crypto interval
markets ask whether the price of Bitcoin, Ethereum or Solana will be above its starting
value at the end of a five- or fifteen-minute interval, and a new one opens as the previous
one closes; their prices follow an external reference, the spot price on a centralised
exchange, with a lag of a few hundred milliseconds \cite{young2026openmarket}.  Sports
markets trade before and during the game, with a recorded start time; political, economic
and cultural questions trade for weeks or months; daily temperature and weather questions
resolve every day.

Resolution is not decided by Polymarket.  When the event is over, a user proposes the
outcome to UMA's optimistic oracle and posts a bond in collateral; the proposal stands if
nobody disputes it within a two-hour challenge window, a wrong or premature proposal
forfeits the bond, and a disputed proposal goes through the oracle's dispute-resolution
process \cite{polymarket2026resolution}.  Resolution is therefore an event that happens
after the end date, sometimes well after it.  The thesis never uses the outcome or the
resolution time as an input: it predicts the price while the market is still open.
\todo{numbers on how many markets exist and how many are active per day, from
\cref{ch:data}}

\subsection{What is known about Polymarket}
\label{sec:pm:related}

Machine-learning work on Polymarket has so far concentrated on \emph{outcome}
forecasting with language models and news \cite{cheng2026polybench,barot2026polyswarm},
on dataset construction \cite{jia2026unlocking,young2026openmarket} and on the economics
of the platform \cite{reichenbach2025exploring,tsang2026anatomy}.
\textcite{young2026openmarket} is the closest in spirit: a logistic model on 43
order-flow features does not beat the Polymarket mid price itself as a predictor of the
outcome of the 15-minute Bitcoin markets.  The present thesis predicts the \emph{price},
not the outcome, over horizons at which the market is still open; the natural references
are therefore the limit-order-book forecasting literature (\cref{sec:bg:lob}) and the
literature on dynamic graphs (\cref{sec:bg:dyn}).  \todo{check whether anything else
came out on price (not outcome) forecasting}

\section{Time series, order books and sequence models}
\label{sec:bg:ts}

\subsection{Concepts}
\label{sec:bg:ts:concepts}

\todo{to write.  Outline:}
\begin{itemize}[nosep]
  \item time series; the order-book and transaction streams are sequences of timed
  \emph{events}, not series sampled on a regular grid;
  \item two ways of reading a stream: a fixed summary of the past handed to a learner
  without memory (\emph{stateless}), or a \emph{state} per market updated at every event
  (\emph{stateful}, recurrent); what each can see, how each is trained, how each handles
  irregular time;
  \item the task: the direction of the change over a horizon $H$, three classes
  (\textsc{down}, \textsc{const}, \textsc{up}); the random walk as the null model.
\end{itemize}

\subsection{Price predictability and limit order books}
\label{sec:bg:lob}

The efficient-market hypothesis in its weak form states that past prices carry no
information about future prices \cite{fama1970efficient}.  At short horizons this is only
approximately true: the stylised facts of high-frequency returns \cite{cont2001empirical}
include a negative autocorrelation at the shortest lags, produced by the bid--ask bounce
and by the discreteness of prices, and in a limit order book the next move of the mid is
related to the order-flow imbalance at the best quotes \cite{cont2014price}.
\todo{deep learning on limit order books, from DeepLOB \cite{zhang2019deeplob} on;
predictability depends on the microstructure more than on the architecture
\cite{briola2024deep}; one model pooled over many assets \cite{sirignano2019universal};
the warning about linear baselines \cite{zeng2023transformers}}

\subsection{Recurrent networks and multi-rate memories}
\label{sec:bg:rnn}

\todo{GRU \cite{cho2014gru} and truncated backpropagation through time; the clockwork RNN
\cite{koutnik2014clockwork} and GRU-D \cite{che2018grud} as the references of the
hierarchical memory and of the continuous-time decay}

\subsection{Time-series foundation models}
\label{sec:bg:tsfm}

\todo{Chronos and Chronos-2 \cite{ansari2024chronos,ansari2025chronos2}; the evidence in
finance, close to a random walk}

\section{Graphs and graph neural networks}
\label{sec:bg:graphs}

\subsection{Concepts}
\label{sec:bg:graphs:concepts}

\todo{to write.  Outline:}
\begin{itemize}[nosep]
  \item graphs; the bipartite graph between wallets and markets; degree, components,
  $k$-cores;
  \item dynamic graphs: a sequence of snapshots or a stream of timed events;
  \item message passing; on a bipartite dynamic graph, from wallets to the markets they
  trade and back.
\end{itemize}

\subsection{Graph neural networks}
\label{sec:bg:gnn}

\todo{message passing \cite{scarselli2009gnn,micheli2009nn4g,gilmer2017mpnn,bacciu2020gentle};
neighbourhood sampling and attention over the neighbours; a GRU as the node update
\cite{li2016ggnn}}

\subsection{Dynamic graphs}
\label{sec:bg:dyn}

\todo{continuous-time dynamic graphs \cite{kazemi2020survey,gravina2024dynamic}; the
memory-based models up to Temporal Graph Networks \cite{rossi2020tgn}; a model should
contain the strong heuristic as a special case \cite{feldman2026navis}}

\section{Gradient boosting and neural tabular learners}
\label{sec:bg:tabular}

\todo{LightGBM \cite{ke2017lightgbm} as the reference on tables of features; how a
neural tabular model has to be trained to compete
\cite{gorishniy2022embeddings,holzmuller2024betterbydefault}; TabM
\cite{gorishniy2025tabm}; logit adjustment under class imbalance
\cite{menon2021logitadjustment}}

\section{Evaluation under imbalance and dependence}
\label{sec:bg:eval}

\todo{balanced accuracy \cite{brodersen2010balanced}; dependent samples and the
hour-block bootstrap \cite{kunsch1989jackknife}; the focal loss \cite{lin2017focal}}
